% =====================================================================
%  Professioni 0b — La comprensione del testo
%  Fonte: trascrizione dell'incontro sulla comprensione del testo dei
%  Dev Days (AIV, "lezione zero"): un incontro online su Zoom di circa
%  un'ora e mezza, il giovedì mattina (inizio previsto alle 10, effettivo
%  alle 10:15 circa), tenuto da Daniele Fusetto, direttore del corso di
%  narrative design di AIV. Da 20 a 26 partecipanti di tutti i corsi del
%  primo anno; slide condivise; tre esercizi a coppie nelle sessioni
%  secondarie di Zoom; domande finali.
%  Pacchetti richiesti nel preambolo del libro:
%    \usepackage[utf8]{inputenc}, \usepackage{booktabs},
%    \usepackage{tikz} con
%    \usetikzlibrary{arrows.meta,positioning,calc,fit}
%
%  Collocazione nel libro: parte "Alcune Professioni nel Mondo
%  Videoludico", subito dopo AIV0_mindset_e_best_practice (3.2): come
%  quello, è un incontro di carattere generale.
%  Tutte le etichette di questo capitolo hanno il prefisso "pt:".
% =====================================================================

\chapter[La comprensione del testo]{La comprensione del testo: termini,
consegne, registro}
\label{pt:cap:comprensione}

Fra gli incontri dei Dev Days (capitolo~\ref{p1:cap:devdays}) ce n'è uno
che non presenta un corso: la ``lezione zero'' sulla \textbf{comprensione
del testo}, obbligatoria per tutti gli studenti del primo anno, tenuta da
\textbf{Daniele Fusetto}, direttore del corso di narrative design (lo stesso
relatore dei capitoli~\ref{p0:cap:mindset}, \ref{p1:cap:devdays}
e~\ref{p3:cap:narrative}). È, dice, ``una \emph{new entry}'': è stata aggiunta
quest'anno perché l'anno scorso ci sono stati dei problemi.

Il motivo lo spiega in apertura. Negli ultimi anni, ``sempre in maniera più
crescente'', si arriva a un punto in cui ``le generazioni nuove, ma anche le
generazioni vecchie, non si riescono a parlare fra di loro'', e nascono
incomprensioni di fondo. Il problema dipende in parte anche dai docenti
(``come vedrete, oggi l'interlocutore è sempre concausa del problema''), in
parte da chi legge. E la scuola italiana, secondo il relatore, negli ultimi
anni ha fatto ``un non lavoro'' su un certo tipo di comprensione: tutti fanno
comprensione del testo alle elementari, alle medie e alle superiori, ma su
testi argomentativi ed esplicativi, non sui testi \textbf{operativi}, quelli
che chiedono di fare qualcosa. L'incontro serve a tre cose:
\begin{itemize}
  \item vivere meglio la classe e avere meno incomprensioni fra compagni;
  \item avere meno incomprensioni con i docenti e i direttori di corso, e capire
        come parlare loro (``sono un po' stufo di mail tipo \emph{ciao bro}
        con lo slang: non voglio le mail formalissime accademiche, però
        nemmeno \emph{ciao bro}'');
  \item capire come si comprende una \textbf{richiesta}: gli studenti avranno
        delle task, faranno progetti in team, e in futuro dovranno comprendere
        le richieste dei datori di lavoro.
\end{itemize}
``Quindi, benvenuti nel mondo dell'incomprensione''.

L'incontro si tiene online su Zoom il \textbf{giovedì mattina}: con ogni
probabilità \textbf{giovedì 8 ottobre 2026}\footnote{Il calendario annunciato
all'apertura della settimana (capitolo~\ref{p1:cap:devdays}) colloca la
comprensione del testo il giovedì alle 10; il relatore dice inoltre che con
narrative design ``iniziamo domani'', e la prima lezione di narrative è
venerdì 9 ottobre (capitolo~\ref{p3:cap:narrative}).}. Doveva cominciare alle
10, ma comincia con un quarto d'ora di ritardo (``altri impegni accademici'');
più avanti il relatore dice che sono ``le 11:23''. I partecipanti sono una
ventina all'inizio (``spero più di lunedì. Beh, vedo le solite facce''), poi
26. È una lezione con \textbf{esercizi}: il relatore apre una dozzina di
sessioni secondarie su Zoom per far lavorare gli studenti a coppie, a tempo.

Avvertenza di lettura: la trascrizione è automatica e, soprattutto durante gli
esercizi, quando gli studenti sono nelle sessioni secondarie e il microfono
principale resta muto, contiene lunghissimi blocchi di testo allucinato dal
software (``il regista, il regista'', ``qual è la risposta che ci può
essere?'', ``la loro dottrina sul mercato del lavoro'', ``il governo di Dalai
Lama'' e simili, ripetuti decine di volte), qui ignorati. Storpia inoltre
molti nomi: AIV diventa ``IV'', ``live'', ``AIF''; il relatore stesso è
chiamato ``Danilo'' e ``professor Furetto''; \emph{curse of knowledge} diventa
``course of knowledge''; Kintsch ``Kinch''; Giles ``Gills''. Dove la
ricostruzione è sicura si usa la forma corretta, altrimenti lo si dice in nota.
Il relatore cita molti studi di linguistica e psicologia cognitiva: dove una
citazione o un dato risultano imprecisi lo si segnala in nota.

Il disclaimer iniziale, ``che non c'è nelle slide'': \textbf{non si possono
risolvere il 100\% dei problemi di comprensione}. L'incomprensione esisterà
sempre, e a volte è positiva, ``un po' come la procrastinazione'': ci hanno
insegnato che procrastinare è sbagliato, ma esistono una procrastinazione
positiva e una negativa. Lo stesso per l'incomprensione: ``a volte
dall'incomprensione nascono giochi fichissimi''. Però ``c'è incomprensione e
incomprensione''.

Il percorso del capitolo, che segue le tre parti della lezione:
\begin{enumerate}
  \item \textbf{lessico e dizionari}: concetto, termine, definizione; falsi
        amici, termini di prassi, prestiti; il dizionario personale e il
        \emph{common ground}; il primo esercizio
        (sezioni~\ref{pt:sec:lessico} e~\ref{pt:sec:es1});
  \item \textbf{lettura e comprensione}: i tre livelli del testo, gli atti
        linguistici, la pertinenza, l'intento del comandante, l'illusione di
        Mosè; il secondo esercizio
        (sezioni~\ref{pt:sec:lettura} e~\ref{pt:sec:es2});
  \item \textbf{registro e contesto}: l'accomodamento, la maledizione della
        conoscenza, il \emph{grounding}, come scrivere una mail a un docente;
        il terzo esercizio (sezioni~\ref{pt:sec:registro}
        e~\ref{pt:sec:es3});
  \item le domande finali: l'inglese, il tripla A o l'indie, il lavoro, gli
        stipendi (sezione~\ref{pt:sec:domande}).
\end{enumerate}

% ---------------------------------------------------------------------
\section{Lessico e dizionari}
\label{pt:sec:lessico}

\subsection{Le norme ISO e il vocabolario controllato}

Si parte da una ``sorta di trivia'' utile. Esistono norme internazionali
standard di carattere tecnico, le \textbf{ISO}; fra queste, la \textbf{ISO
1087}\footnote{ISO 1087, \emph{Terminology work and terminology science —
Vocabulary}.} definisce il vocabolario fondamentale della terminologia, in
inglese e a volte anche in francese. ``Sono molto noiosi, ve lo dico''. Ma
sono vocabolari fondamentali per accordarsi quando si fa qualcosa di molto
importante, ``vitale'': se si è scienziati in Antartide, ``lasciare la porta
aperta non è come a casa, che tua madre ti sgrida: lì è la vita o la morte''.

Un dizionario ISO fornisce termini e concetti chiave, li classifica e
formalizza per gerarchie, e regola la distinzione fra definizioni
\emph{intensionali} ed \emph{estensionali}\footnote{Il relatore le spiega come
``termini che usiamo in maniera specifica per un dato motivo, o che per
estensione usiamo per qualche altro motivo''. Nella ISO 1087 la distinzione è
un'altra: una definizione \emph{intensionale} descrive un concetto attraverso
le sue caratteristiche (il genere prossimo e le differenze), una definizione
\emph{estensionale} lo descrive elencando gli oggetti o i concetti subordinati
che vi rientrano.}. Nello sviluppo di videogiochi l'ISO non si usa: è ``troppo
complesso, troppo statico, troppo inamidato''. Ma fa capire che esiste un modo
di organizzare il lavoro di centinaia di persone sul piano terminologico:
\textbf{fare un dizionario}.

\subsection{Concetto, termine, definizione}

Dalla ISO 1087 il relatore prende tre concetti, per mostrare che anche questi
hanno una loro differenza. Una parola della lingua comune può avere un
significato; un \textbf{termine tecnico} ha invece un \textbf{dominio} ben
preciso: se nella game industry si parla di \emph{build}, si indica una cosa
completamente diversa da quella che indica la stessa parola nel mondo
dell'acciaio.
\begin{itemize}
  \item Il \textbf{concetto} è l'unità di conoscenza, ciò che si intende. In
        semiotica si distinguono il \emph{referente} (la cosa a cui ci si
        riferisce, il cane reale), il \emph{significato} (il concetto di cane)
        e il \emph{significante} (la parola ``cane''); ``a volte ci sono delle
        crasi''.
  \item Il \textbf{termine} è l'etichetta linguistica assegnata a quel concetto
        in un \textbf{dominio specifico}. Un veterinario, a un convegno di
        biologi, non chiamerà il cane ``cane'', ma \emph{Canis lupus
        familiaris}: non ovunque, ma in quel contesto.
  \item La \textbf{definizione} è la descrizione che delimita il concetto
        rispetto a quelli vicini.
\end{itemize}
Chiedere ``mi date la definizione che voi intendete per quel termine?''
permette di risolvere il problema dei \textbf{falsi amici}. Se il relatore parla
di \emph{prototipo} con i suoi amici architetti, a loro vengono in mente i
modellini con gli stecchini dell'università; per uno sviluppatore il prototipo
è una \emph{build} di gioco fatta per valutare alcune meccaniche con il
\emph{playtest} (``un prototipo senza playtest di solito è inutile''). E le
build stesse si distinguono: milestone, \emph{branch build}, \emph{main
build}; e i narrative chiamano le build in modo diverso dai programmatori.
Spesso c'è addirittura un significato implicito, recondito, nel termine, che
si imparerà durante il corso.

\subsection{La maledizione della conoscenza, i falsi amici, i termini di prassi}

Soprattutto all'inizio, attenzione: i docenti useranno parole che daranno per
scontate. È un bias che ha un nome, la \emph{curse of knowledge}
(sezione~\ref{pt:sec:registro}). Il relatore, che lavora nell'industria, dà
per scontato che chi gli sta intorno capisca quello che dice; cerca di
semplificare il più possibile nelle prime lezioni (che fa in inglese), ma gli
scappano termini tecnici. Tre trappole:
\begin{itemize}
  \item i \textbf{falsi amici}, che esistono fra l'inglese e l'italiano ma anche
        nella lingua tecnica. L'esempio ``canonico'' è l'\textbf{MDA}, un
        framework di game design, con le sue tre parole: meccaniche, dinamiche
        ed \emph{estetiche}. Che cos'è l'estetica? ``L'aspetto estetico,
        no?''. Viene da pensare che le estetiche di un gioco siano l'aspetto
        grafico: è sbagliato, \textbf{le estetiche sono le emozioni che il
        gioco crea}\footnote{Nel framework MDA (Hunicke, LeBlanc e Zubek, 2004)
        le \emph{aesthetics} sono le risposte emotive desiderate nel giocatore.
        Si veda anche la parte di Game Design di questo libro.}. Il relatore ha
        visto questo errore ``tantissimo'' nell'industria, ``con problematiche
        enormi: quando capisci male una teoria e la applichi male, poi ti
        ritrovi con delle ciofeche'';
  \item i \textbf{termini di prassi}: termini che non hanno doppi significati, ma
        che si dà per scontato di conoscere. ``Io dico \emph{pitch} e voi dite:
        ah, beh, so cos'è un pitch''. Poi si scopre che il pitch del relatore è
        diverso da quello di Filippo, di Eleonora, di Gabriele, di Emanuele: e
        se si è tutti nello stesso team, è il caos;
  \item i \textbf{prestiti} dall'inglese o da altre lingue, ``veramente
        difficili''.
\end{itemize}
Il relatore mostra uno studio linguistico che classifica i tipi di falsi
amici\footnote{L'immagine proviene da uno studio comparativo fra inglese,
spagnolo, francese e forse giapponese; la trascrizione non permette di
identificarlo, e gli esempi sono in parte storpiati.}: i \emph{true friends},
parole identiche con lo stesso significato (\emph{adventure},
\emph{aventure}); i falsi amici totali, identici nella scrittura e nella
pronuncia ma non nel significato; i falsi amici parziali, come l'inglese
\emph{ancient} e l'italiano ``anziano'', ``in cui uno pensa di arrivarci, ma
non ci arriva davvero''; e i parziali ``di secondo tipo'', come l'inglese
\emph{demand} e il francese \emph{demander} (``chiedere'').

\subsection{Non esiste un dizionario del game development}

Non esiste un dizionario di game design, né di game development in generale,
``e questo vale per tutti, anche per i programmatori, anche per i 3D artist''.
È anche il motivo per cui, a differenza delle superiori o dell'università, in
AIV non ci sono libri di testo: si consigliano libri, ma non c'è ``il'' libro.
In altre discipline ci sono testi talmente canonici da essere chiamati con i
nomi degli autori: in narratologia c'è ``l'Onega-Landa''\footnote{Susana Onega
e José Ángel García Landa (a cura di), \emph{Narratology: An Introduction}
(1996). Nella trascrizione ``L'Oniga Landa''.}, una raccolta di saggi con
un'introduzione che è una pietra miliare. Nel videogioco no, per due motivi:
\begin{itemize}
  \item la disciplina è \textbf{giovane e cambia}: un manuale di JavaScript
        comprato oggi fra sei mesi, otto mesi, un anno è già vecchio; un
        manuale di ChatGPT non ha più senso appena esce la versione successiva;
  \item ogni azienda ha il suo software e il suo \emph{know-how}. Le aziende
        ``fatte bene'' hanno un loro \textbf{dizionario di produzione}, un
        dizionario di area, un loro slang e un loro gergo.
\end{itemize}

\subsection{Il dizionario personale}

Per questo il relatore invita a farsi un \textbf{dizionario personale}:
prendere appunti a lezione, segnarsi le parole che non si conoscono, chiederne
il significato al docente, e soprattutto assegnare loro dei \textbf{sinonimi},
``perché il modo migliore per imparare i termini non è dare una definizione, ma
mettere sinonimi''. Chiedere ``qual è un sinonimo di questo?'', o mettere i
propri; poi chiedere anche la definizione o una delucidazione, ``non è un
problema''. Nel game design c'è addirittura uno strumento, il
\textbf{production dictionary}, per cui tutto il team si mette a un tavolo e
concorda i termini tecnici di base.

\subsection{Il common ground}

Un altro metodo per disambiguare mail, dialoghi e conversazioni è trovare il
\textbf{common ground}: l'insieme di conoscenze comuni fra due o più persone,
``il modo migliore per partire in un progetto''. Di solito è fatto di parole o
di esperienze che, nella propria esperienza individuale, sono condivise. Il
relatore lo fa sempre, anche in consulenza: nella chiacchierata iniziale di un
team, quando nella stanza si parla del più e del meno, chiede ``qual è l'ultimo
gioco che avete giocato? Avete giocato \emph{Resonant}?''\footnote{Con ogni
probabilità \emph{Control Resonant}, di cui il relatore parla anche nel
capitolo~\ref{p3:cap:narrative}.}. È una tecnica di common ground: valuta se il
creative director ha avuto le sue stesse esperienze, se il game director ha
giocato i suoi stessi giochi, se certe reference di meme (``\emph{this is
fine}'') vengono colte. Così si parte col piede giusto.

\subsection{I documenti sono per tutti}

Nei team grandi e nei documenti, il common ground va sostituito con degli
\textbf{strumenti}: documenti di contesto, il \emph{syllabus} del corso (ogni
corso di AIV ne ha uno), glossari interni del team, documenti di progettazione.
E qui una richiesta ``ai miei amici game designer'': \textbf{fate i documenti
per tutti, non solo per voi}. ``Non fate documenti di game design scritti in
egiziano'': il GDD non è il documento \emph{del game designer}, è il documento
\emph{del game design}, e il gioco è di tutti. Se è pieno di parole astruse e
termini tecnici che il 3D artist e il programmatore non conoscono, ma che
devono usare, ``ciaone, avrete dei grossi problemi''. Chiedersi sempre: questo
documento è solo per uso interno, o devo spiegare qualcosa a qualcuno? Vale
anche per i programmatori, ``che però di solito questo errore lo fanno un po'
meno''.

Altri consigli:
\begin{itemize}
  \item confrontare come lo stesso termine compare in \textbf{testi diversi}
        (gli usi multipli);
  \item fare \textbf{ricerche online};
  \item fare sempre degli \textbf{esempi}, e quando si nota un'incomprensione
        chiedere: ``ma tu stai parlando del pitch nel senso del pitch deck che
        abbiamo fatto la settimana scorsa?''. Molto spesso le incomprensioni
        nascono da qui.
\end{itemize}

% ---------------------------------------------------------------------
\section{Primo esercizio: le definizioni}
\label{pt:sec:es1}

Il primo esercizio ha una fase individuale e una a coppie, nelle sessioni
secondarie di Zoom. La consegna:
\begin{enumerate}
  \item individualmente, scrivere \textbf{una riga di definizione} per ciascuno
        di sette termini: \emph{prototipo}, \emph{iterazione}, \emph{build},
        \emph{scope}, \emph{milestone}, \emph{deliverable}, \emph{polish}. Se non
        si sa, non si cerca online: si scrive la definizione che sembra più
        vicina al significato, o niente;
  \item a coppie, confrontare le risposte e gli scarti: individuare i falsi
        amici, i termini di prassi, i prestiti ``spostati'', che sono
        soggettivi e relativi alla coppia;
  \item dopo dieci minuti, in plenaria, sentire alcune coppie: chi aveva falsi
        amici, chi definizioni uguali, chi la definizione ``più astrusa e
        assurda''.
\end{enumerate}
L'ingresso nelle sessioni è a sua volta una piccola lezione di comprensione:
alcuni studenti entrano prima che il relatore lo dica (``non vi ho ancora detto
di entrare, state qui, calmi''; ``ecco, meno male che abbiamo fatto la
comprensione del testo''; ``vedete che serviva?''), altri non trovano il
comando (``in basso a destra, tre pallini, \emph{sessioni secondarie}'').

\subsection{La discussione}

Più che discrepanze, uno studente si è accorto di non conoscere proprio alcuni
termini; altri hanno tirato a indovinare. I termini discussi
(tabella~\ref{pt:tab:termini}):
\begin{itemize}
  \item \textbf{deliverable}: la parola che più ha messo in difficoltà. Uno
        studente aveva scritto ``prodotto consegnabile'', traducendo
        letteralmente; il compagno aveva capito ``derivabile'', ``una cosa che
        deriva da un concetto più generale, un'istanza specifica'' (``quello è
        \emph{derivable}''). Nella game industry il deliverable è \textbf{ciò
        che ci si aspetta venga consegnato}: ``qual è il deliverable che io mi
        aspetto da voi?'', o anche il deliverable che si promette di consegnare
        a una certa data;
  \item \textbf{iterazione}: ``che italiano è?'', scherza il relatore. Uno
        studente risponde ``la ripetizione di un'azione''; un altro aveva
        pensato a \emph{iter} come percorso. È vicino alla ripetizione:
        iterare significa \textbf{rifare un ciclo} di creazione della build,
        playtest, prototipazione, feedback. Sviluppare in modo iterativo è
        tipico dell'industria (si veda il ``fusillo'' del
        capitolo~\ref{p2:cap:gamedesign});
  \item \textbf{scope}: ``oh, questa è molto bella''. Lo scope di un progetto è
        il suo obiettivo in termini di \textbf{quantità}, di impegno, di
        \emph{effort}: ``lo scope di questo progetto non è fare un open world,
        ma un gioco con tante aree separate che sembra un open world''. Viene
        dall'inglese, ``ambito, portata''; anche ``ambizione''. Si usa molto il
        concetto di \textbf{overscope}: produrre un gioco oltre l'effort che ci
        si può mettere. Lo scope è legato alle milestone: all'inizio si
        definiscono i propositi in termini \textbf{concreti} (``il gioco avrà 7
        livelli, 3 personaggi giocabili, 7 meccaniche di combattimento''); è
        questa la differenza fra scopo e obiettivo. Le \textbf{milestone} sono
        invece gli step con cui si svilupperà il progetto.
\end{itemize}
Filippo aveva inteso lo scope in senso commerciale, come obiettivo di vendita o
di clientela. È possibile, risponde il relatore, ma nel marketing e nelle
vendite si usano termini diversi: per le persone \emph{target} o
\emph{segmento di mercato}; nel digital e nei social lo scope si misura con i
\textbf{KPI}\footnote{Nella trascrizione ``key point indicator'': la sigla
significa \emph{key performance indicator}.} (``con questa campagna vogliamo
passare da 500 a 1.000 follower''). Termini che forse non serviranno, ma da
ricordare se si parla con qualcuno del marketing.

\begin{table}[htbp]
\centering
\small
\begin{tabular}{@{}lp{9.6cm}@{}}
\toprule
\textbf{Termine} & \textbf{Significato nella game industry (dalla lezione)} \\
\midrule
prototipo & build di gioco fatta per valutare alcune meccaniche con il
playtest (non il modellino dell'architetto) \\
iterazione & ripetere il ciclo build, playtest, prototipazione, feedback \\
build & versione eseguibile del gioco; ci sono milestone build, branch build,
main build \\
scope & portata del progetto in termini concreti di quantità ed effort;
overscope: oltre l'effort sostenibile \\
milestone & gli step con cui si sviluppa il progetto \\
deliverable & ciò che ci si aspetta (o si promette) venga consegnato \\
polish & (non discusso a voce; di norma, la rifinitura finale) \\
\bottomrule
\end{tabular}
\caption{I sette termini del primo esercizio. La definizione di
\emph{polish} non è stata discussa in plenaria.}
\label{pt:tab:termini}
\end{table}

% ---------------------------------------------------------------------
\section{Lettura e comprensione}
\label{pt:sec:lettura}

\subsection{I tre livelli del testo}

Walter Kintsch\footnote{Nella trascrizione ``Walter Kinch''. Il modello è
quello di Kintsch e van Dijk, poi sviluppato da Kintsch nel modello di
costruzione e integrazione.} ha suddiviso il testo che un lettore legge in
\textbf{tre livelli} (figura~\ref{pt:fig:kintsch}):
\begin{itemize}
  \item la \textbf{superficie} (\emph{surface form}): le parole e la sintassi,
        o il disegno; è ciò che si dimentica più in fretta;
  \item la \textbf{base testuale} (\emph{textbase}): le proposizioni, ciò che
        il testo afferma;
  \item il \textbf{modello della situazione} (\emph{situation model}): la
        rappresentazione mentale che ci si fa, anche inconsciamente.
\end{itemize}
L'esempio della slide è una rana che mangia un insetto. Il relatore nota che la
slide sembra invertita (la superficie al posto del modello della situazione) e
poi capisce perché: si può andare in entrambi i sensi. Se si dà un disegno, la
superficie è il disegno; se si dà la frase \emph{the frog ate the bug}, la
superficie è la frase. La base testuale resta la proposizione (mangiare, rana,
insetto), e il modello della situazione è l'immagine di una rana che mangia un
insetto. Non tutti ``vedono'' le immagini (esiste l'\emph{afantasia}), ma la
mente fa comunque le connessioni fra la parola \emph{frog} e l'idea di rana.

\begin{figure}[htbp]
\centering
\begin{tikzpicture}[>=Stealth, font=\small, node distance=0.6cm,
  box/.style={draw, rounded corners, align=center, text width=3.4cm,
              minimum height=1.2cm, font=\scriptsize}]
  \node[box] (s) {\textbf{superficie}\\parole e sintassi\\\emph{``The frog ate the bug''}};
  \node[box, right=of s] (t) {\textbf{base testuale}\\proposizioni\\MANGIARE(rana, insetto)};
  \node[box, right=of t, fill=gray!15] (m) {\textbf{modello della situazione}\\l'immagine mentale:\\quale insetto? che rana?};
  \draw[->] (s) -- (t); \draw[->] (t) -- (m);
  \node[below=0.3cm of m, font=\scriptsize\itshape, align=center, text width=3.6cm]
    {qui entra ciò che il lettore aggiunge di suo};
\end{tikzpicture}
\caption{I tre livelli del testo secondo Kintsch, con l'esempio della slide.
Ricostruzione schematica.}
\label{pt:fig:kintsch}
\end{figure}

\paragraph{Quale insetto?} Il relatore chiede di scrivere in chat il primo
insetto che viene in mente sentendo \emph{the frog ate the bug}: scarafaggio,
mosca, zanzara, cimice. Molti dicono mosca, ma non tutti: la parola
\emph{insetto} è \textbf{vaga}, e il modello della situazione è diverso per
ognuno. Lo stesso con la rana: di che colore è? Verde, o magari, se la sera
prima si è visto un documentario sul Sud America, ``quella nera e blu che se la
leccate fate dei viaggioni cosmici''\footnote{Le rane ``freccia'' (famiglia
Dendrobatidae) del Sud America, spesso nere e blu, sono \emph{velenose}, non
allucinogene: la leggenda del ``leccare'' riguarda semmai il rospo del deserto
di Sonora (\emph{Incilius alvarius}), le cui secrezioni contengono sostanze
psicoattive.}. Dipende molto dal contesto.

\paragraph{Il problema dei testi operativi.} Nei testi operativi si agisce sul
\textbf{modello della situazione}, perché bisogna decidere che cosa fare; ed è
un problema enorme, perché lo si costruisce a partire da un testo che può essere
molto ambiguo. Due persone possono leggere la stessa superficie, estrarne basi
testuali diverse e costruire due compiti completamente diversi. \textbf{L'errore
sta in ciò che il lettore aggiunge di suo, o non aggiunge}. Se il relatore dice
``per domani fatemi un documento con tutte le parole che ho detto e che non
avete capito'', qualcuno userà Google Docs, qualcuno un .docx, qualcuno
OpenOffice, qualcuno manderà un file di testo (``probabilmente è un
programmatore''), un artista magari un JPEG. Più si è precisi, meglio è; ma
bisogna ricordarsi anche dell'interlocutore.

\subsection{Gli atti linguistici}

Come si interpreta un comando? Con la teoria degli \textbf{atti linguistici}
(Austin, Searle): spesso \textbf{dire è anche fare}. In ogni enunciato si
distinguono (tabella~\ref{pt:tab:atti}):
\begin{itemize}
  \item l'\textbf{atto locutorio}: ciò che viene detto;
  \item l'\textbf{atto illocutorio}: l'azione che si compie dicendo;
  \item l'\textbf{atto perlocutorio}: l'effetto prodotto.
\end{itemize}
Quando il relatore ha detto ``adesso apro le aule'', non ha detto ``entrate'':
implicitamente chiedeva, o ordinava, di entrare, e l'effetto è chi entra e chi
no. La distinzione è utilissima per classificare gli errori: a volte l'errore
è in ciò che viene detto, altre in ciò che viene compreso, altre ancora
nell'effetto prodotto. Quando l'effetto è del tutto imprevedibile, può esserci
un disturbo nell'interlocutore: se si chiede ``mi passi il sale?'' e arriva lo
zucchero, probabilmente l'altro era distratto: non ha frainteso il comando, la
sua risposta è stata disturbata dal \textbf{rumore}. È molto comune, e insegna
che bisogna \textbf{ascoltare e non fare multitasking}, ``come vi abbiamo detto
lunedì con Andrea'' (capitolo~\ref{p1:cap:devdays}).

Le slide riportano anche la classificazione degli atti (dichiarativi,
rappresentativi, commissivi, direttivi, espressivi: ``se volete approfondire,
andate ad approfondire'') e la distinzione fra tre ``forze'': la forza
locutoria, ``quanto è forte il messaggio''; quella illocutoria, ``quanto è forte
l'intenzione''; quella perlocutoria, ``quanto è forte l'effetto inteso
sull'ascoltatore''\footnote{Nella teoria classica il concetto centrale è la
\emph{forza illocutoria} (il tipo di atto: ordine, promessa, domanda\dots);
parlare di ``forza locutoria'' e ``forza perlocutoria'' come gradi di intensità
è una semplificazione didattica del relatore.}. Quest'ultima è fondamentale
anche nei videogiochi, perché ``i videogiochi sono comunicazione
differita''\footnote{Nella trascrizione ``indifferita''.}: se si dice male al
giocatore ``usa l'ascia per spaccare quella porta'', il giocatore si lancia
direttamente contro la porta, come nel meme; o come nel meme di Peter Griffin
``quella volta che mi sono dimenticato come ci si siede''.

\begin{table}[htbp]
\centering
\small
\begin{tabular}{@{}lp{4.3cm}p{5.2cm}@{}}
\toprule
\textbf{Atto} & \textbf{Che cos'è} & \textbf{Esempio della lezione} \\
\midrule
locutorio & ciò che viene detto & ``Adesso apro le aule'' \\
illocutorio & l'azione compiuta dicendo & la richiesta implicita: entrate \\
perlocutorio & l'effetto prodotto & chi entra e chi no; chi entra prima del
tempo \\
\bottomrule
\end{tabular}
\caption{I tre livelli dell'atto linguistico, con l'esempio delle sessioni
secondarie.}
\label{pt:tab:atti}
\end{table}

\subsection{La pertinenza}

Per interpretare un comando il relatore prende due teorie. La prima è la
\textbf{teoria della pertinenza}\footnote{La teoria della pertinenza
(\emph{relevance theory}) è di Dan Sperber e Deirdre Wilson (1986).}, che
aggiunge una cosa: il lettore \textbf{non decodifica il testo parola per
parola}. Si legge per contesto: è anche il motivo per cui, imparando una
lingua, si parte dai vocaboli, poi si legge e dal contesto se ne capiscono
altri, che vengono integrati nel proprio ``dizionario mentale''. L'errore tipico
è interpretare male la pertinenza: sentire una parola usata in un contesto
sbagliato e, per l'\textbf{effetto Einstellung} (sezione~\ref{p0:sec:limiti}),
continuare a usarla nei contesti sbagliati; correggersi poi è molto difficile.

Per questo uno dei problemi principali non è la lingua imparata a scuola, ma il
\textbf{patois}: il dialetto familiare o della propria comunità. In casa, o fra
amici, si usano parole che fuori sono usate diversamente, battute che altri non
colgono: ``non a tutti fa ridere il numero 67, ma se lo dico in una classe di
bambini delle elementari impazziscono'' (Sara spiega: per il meme \emph{six
seven}).

\subsection{L'intento del comandante}

La seconda teoria viene dalla dottrina militare: il \textbf{mission
command}\footnote{Derivato dalla \emph{Auftragstaktik} prussiana; il concetto
chiave è il \emph{commander's intent}, l'intento del comandante, che permette a
chi esegue di adattarsi senza perdere lo scopo.}, ``tranquilli, non vi voglio
trasformare in soldati''. È la disciplina, insegnata ai generali, che insegna a
dare ordini e scopi chiari, soprattutto in situazioni di grande stress. La
slide la confronta con il \emph{command and control}; si usano in momenti
diversi. È interessante notare che nella dottrina militare, al netto
dell'addestramento (i film sui Marines, ``ti voglio schiacciare''), c'è
tantissimo lavoro di \textbf{ispirazione} e di \emph{enabling}: se un team di
Navy SEAL va in missione segreta, sotto pressione, in una parte del mondo da cui
sa di poter non tornare, e il comandante dà gli ordini dicendo ``boh, ragazzi,
non lo so, sperate di tornare'', nessuno avrà il coraggio di partire. Il
\emph{command and control} è il momento della pressione dall'alto, in cui si
evita che chi sta sotto nella gerarchia ``salga''; il \emph{mission command} è
quello in cui il direttore fa l'opposto: abilita le azioni, spinge, supporta,
assiste anche emotivamente.

La \textbf{sintesi operativa}: davanti a una richiesta chiedersi sempre
\textbf{che cosa mi chiede di fare}, \textbf{perché me lo chiede}, e qual è la
pertinenza. Poi c'è tutto il \textbf{non detto}, che il relatore salta per
ragioni di tempo\footnote{Nella trascrizione ``qui c'è, vabbè, il polgrisma'':
con ogni probabilità il riferimento è a Grice e alle sue massime
conversazionali, cioè alla teoria dell'implicatura (il non detto).}.

\subsection{L'illusione di Mosè}

``Quanti animali di ogni specie portò Mosè sull'arca? Rispondetemi rapidamente
in chat''. In chat: ``due per specie, due per specie, due per specie''; qualcuno
``uno per specie''. ``Peccato che non era Mosè, ma Noè''. È l'\textbf{illusione
di Mosè}, studiata da Erickson e Mattson\footnote{Thomas D. Erickson e Mark E.
Mattson, \emph{From words to meaning: A semantic illusion} (1981).}: nei
comandi rapidi, veloci, sotto pressione, non si nota un dettaglio errato ma
\textbf{coerente con il contesto}, che passa inosservato. È stata studiata
negli Stati Uniti, dove i due nomi (\emph{Moses}, \emph{Noah}) non si somigliano
molto; in italiano l'inganno è anche più forte, perché ``Mosè'' e ``Noè'' sono
entrambi tronchi e accentati.

Perché succede? Dall'input semantico si costruisce una struttura coerente (ha
fatto l'arca, ha preso due animali); quando poi bisogna rispondere in fretta, la
mente si sofferma sul dettaglio importante e non su quello ``inferiore''. È
anche un problema di come sono scritte le domande: chiedendo ``ma proprio
Mosè, la persona di Mosè, quanti animali ha portato sull'arca?'', sottolineando
la parola, la gente se ne accorge.

La lezione per i testi operativi: la parte che si dimentica di più è quella
\textbf{nuova e insolita} rispetto alla parte normale. Se si ricevono tre task
uguali di fila, le parti uguali prendono il sopravvento e fanno sbagliare anche
sulle parti nuove. Se la task è sempre ``per domani tre mood board su un
personaggio fatte con Miro'', e la quarta volta il docente chiede di farle con
Milanote, è facile darlo per scontato e farle comunque su Miro.

% ---------------------------------------------------------------------
\section{Secondo esercizio: smontare una consegna}
\label{pt:sec:es2}

Il secondo esercizio si fa nelle stesse coppie. Il testo della consegna,
ipotetica, è sulla slide:
\begin{quote}
Create una mood board per il main character visto a lezione usando un
software di lavagna digitale come Miro o Milanote. Potete suddividere la mood
board se diventa eccessivamente grande, ma non è obbligatorio. La mood board
verrà valutata sulla base della vicinanza al personaggio visto a lezione, la
completezza delle informazioni, la chiarezza espositiva e l'impegno
profuso.\footnote{Ricostruzione dalla lettura a voce del relatore; la
trascrizione scrive ``completazza'' e ``Milanot''.}
\end{quote}
Il relatore non spiega niente della consegna: ``se ci sono cose che non sapete,
andatele a cercare'' (anche cos'è una mood board: ``fa parte del problema della
consegna''). Bisogna \textbf{riscriverla} o parafrasarla, senza copiarne le
frasi (le parole sì, ``per forza''), in quattro colonne
(tabella~\ref{pt:tab:colonne}); si può fare anche a mano, su un quaderno.

\begin{table}[htbp]
\centering
\small
\begin{tabular}{@{}p{2.8cm}p{3.2cm}p{3.2cm}p{3cm}@{}}
\toprule
\textbf{1. Richiesta} & \textbf{2. Livelli} & \textbf{3. Perimetro} &
\textbf{4. Implicito} \\
\midrule
che cosa va fatto e consegnato, e in che forma & che cosa è obbligatorio,
raccomandato, facoltativo & quantità, eccezioni, scadenze e date (se ci sono)
& ciò che è dato per scontato o che manca \\
\addlinespace
\emph{una mood board del main character, su una lavagna digitale} & \emph{la
suddivisione è facoltativa; Miro o Milanote sono esempi, non obblighi} &
\emph{nessuna data né ora di scadenza} & \emph{che cos'è una mood board; dove e
come consegnarla} \\
\bottomrule
\end{tabular}
\caption{Le quattro colonne del secondo esercizio (la slide prevedeva anche una
quinta riga, una regola) e, in corsivo, le risposte emerse in plenaria.}
\label{pt:tab:colonne}
\end{table}

\subsection{La discussione}

``Com'è andata?''. Meglio di prima, rispondono; e mentre scrivevano, rileggendo
la regola, gli studenti si accorgevano di nuove cose. Le domande del relatore:
\begin{itemize}
  \item che cosa manca in maniera emblematica, e dovrebbe esserci in
        \textbf{tutte} le consegne? \textbf{La data di scadenza}, e l'ora:
        questa consegna non ce l'ha;
  \item qual è una parte facoltativa? \textbf{Suddividere la mood board};
  \item su che cosa si viene valutati? La vicinanza al personaggio,
        la completezza delle informazioni, la chiarezza dell'esposizione,
        l'impegno;
  \item che cos'è una mood board? Sara: ``un collage di foto, di solito anche
        con indicazioni disegnate sopra'', che si dà agli artisti, o ad altri,
        per far capire che cosa si intende.
\end{itemize}
Ci sono poi impliciti più complessi. ``Come Miro e Milanote'' \textbf{non è un
obbligo}: si potrebbero usare anche i disegni di Google Docs o delle slide,
perché non c'è scritto ``usate esclusivamente''. Se il docente voleva proprio
quei due software, non ha messo l'obbligo, ed è un problema: in caso di dubbio
\textbf{si chiede} (``alzate la mano: ma quello è obbligatorio?''). Anche perché
spesso le cose si fanno in velocità. ``Questo non ci protegge da delle scuse,
attenzione'': a volte sono i docenti a dare task non chiarissime, altre volte
sono gli studenti a non capire.

% ---------------------------------------------------------------------
\section{Registro e contesto}
\label{pt:sec:registro}

``Cerchiamo di correre, sono le 11:23''.

\subsection{Il registro linguistico}

In linguistica il \textbf{registro} è la \textbf{variazione diafasica}: lo
stesso parlante usa forme diverse a seconda della situazione, del rapporto con
l'interlocutore e dello scopo. È stato studiato anche all'interno di uno stesso
contesto: nei tribunali americani, l'avvocato usa un registro diverso con i
giurati, con un testimone, con il pubblico, con il giudice. Lo stesso contenuto
può essere visto in modi diversi a seconda del registro. E leggere il registro
aiuta sempre a capire il contesto: chi scrive, che ruolo ha, quanto è
vincolante la richiesta, come ci si aspetta che si risponda.

\subsection{L'accomodamento}

Nel dialogo (meno nella scrittura) aiuta il concetto di
\textbf{accomodamento}, studiato da Giles\footnote{Howard Giles, teoria
dell'accomodamento comunicativo (\emph{communication accommodation theory}).
Nella trascrizione ``Gills''.} e da altri studiosi: chi parla tende ad avvicinarsi al
modo di parlare degli altri, soprattutto di chi sta più in alto nella gerarchia.
È l'\textbf{accomodamento di convergenza}: in un team, ``di solito il leader è
la persona di cui gli altri copiano il modo di parlare''. C'è anche
l'\textbf{accomodamento di divergenza}: usare di proposito un modo di parlare
diverso per un effetto preciso, per esempio per mostrare che non si è
d'accordo. Al relatore succede spesso nei playtest: le persone danno feedback
entusiasti (``che bello, che figo''), e chi non è d'accordo dà il suo feedback in
modo sommesso (``boh, non so, non mi sono trovato molto''), di solito per
evitare il conflitto; chi invece diverge per creare conflitto interrompe e dice
con forza ``no, io non sono d'accordo''.

\subsection{La maledizione della conoscenza}

Un esempio celebre viene saltato per mancanza di tempo\footnote{Nella
trascrizione ``l'esempio di Isabel Newton'': con ogni probabilità
l'esperimento di Elizabeth Newton (1990), in cui chi tamburellava il ritmo di
una canzone nota sovrastimava enormemente la probabilità che gli ascoltatori
la riconoscessero.}. Camerer, Loewenstein e Weber hanno studiato la
\textbf{maledizione della conoscenza} (\emph{curse of knowledge})\footnote{Colin
Camerer, George Loewenstein e Martin Weber (1989). Nella trascrizione ``course
of knowledge''.}, a cui tutti sono soggetti, ``soprattutto noi docenti''. È il
motivo per cui il relatore, alle superiori, prendeva alternativamente 4 e 8 in
matematica: la professoressa spiegava dal punto di vista delle proprie
conoscenze e dava per scontate cose che gli studenti non sapevano. La vignetta
della slide: il poliziotto esperto che spiega la strada al passante
(``a destra, poi a sinistra, poi giri qui, poi di là'') e il passante che non si
ricorda niente.

Succederà anche a lezione: il relatore stesso ``pecca'', soprattutto al terzo
anno, perché certe cose sono per lui talmente scontate che è difficile capire
come spiegarle. Per questo bisogna \textbf{chiedere sempre chiarimenti}. Si dice
spesso che chi sa spiegare in modo più elegante e sintetico è chi sa meglio una
cosa; ``all'inizio ero d'accordo, adesso a volte mi fa incazzare''. Non è
sbagliato: chi dà una definizione accurata in poco tempo di solito conosce bene
l'argomento. Ma anche chi fa fatica a spiegarsi può farlo per questo bias: sa
talmente tante cose che fatica a ricostruire la prospettiva di chi non sa.

\subsection{Voce e scrittura}

È stato dimostrato che parlarsi a voce funziona meglio delle mail: avere una
spiegazione vocale della task, oltre al testo. Per questo il relatore darà
sempre il testo della task e lo spiegherà anche a voce. Kruger, Epley e
colleghi\footnote{Justin Kruger, Nicholas Epley e colleghi, \emph{Egocentrism
over e-mail: Can we communicate as well as we think?} (2005). Kruger è lo
stesso dell'effetto Dunning-Kruger, come dice il relatore; la trascrizione
scrive ``Apple''.} hanno chiesto a una serie di persone di scrivere mail
sarcastiche e ironiche: molti erano troppo sicuri di farsi capire, e i
destinatari non le capivano. Quando si scrive, si dà per scontato di essere
chiari. È anche il motivo per cui tutti i regolamenti di gioco che il relatore
legge al primo anno hanno il problema dell'\textbf{egocentrismo narrativo}:
``io so molto di più di te che leggi''.

\subsection{Il grounding}

Nella comunicazione a voce, un modo per evitare questi problemi è il
\textbf{grounding}: tutte le tecniche per verificare se ci si sta spiegando
bene. ``Tutto bene? Mi date un feedback? State capendo?'' è grounding; anche
domandare, cercare informazioni, cercare un accordo. La tecnica di grounding
più fondamentale e utile nel business è \textbf{ricapitolare}: ``posso
ricapitolare? Quindi dobbiamo consegnarvi il materiale in PDF entro il 27
ottobre a questa mail''; ``no, dovete farlo entro quest'altra data, e non in
PDF''. E magari l'errore era proprio nel cliente, che si era spiegato male.
Quindi: sempre domandare.

E ricordarsi che \textbf{il tono percepito in una mail è una proiezione di chi
legge}. ``A volte vi mando una mail rapida, \emph{ciao, grazie, arrivederci}, e
voi: oh mio Dio, com'è stato freddo il prof. Non sono stato freddo, ho poco
tempo, ragà. Non leggeteci Baci Perugina''.

\subsection{La mazza e la palla}

Un altro esperimento: ``una mazza e una palla costano in tutto un euro e dieci.
La mazza costa un euro in più della palla. Quanto costa la palla?''. Qualcuno
in chat dà subito la risposta giusta: \textbf{5 centesimi}, non 10. Si è spinti
a rispondere 10 centesimi; ma allora la mazza costerebbe 1,10 e il totale
sarebbe 1,20. La mazza costa 1,05 e la palla 0,05\footnote{A voce il relatore
dice ``la mazza costa 1,5 euro'': è un lapsus per 1,05.}. È il \emph{Cognitive
Reflection Test} di Frederick\footnote{Shane Frederick, \emph{Cognitive
Reflection and Decision Making} (2005).}: mostra che la mente, se non è molto
attenta, elabora le informazioni in modo troppo rapido.

Sulla slide c'è anche uno schema ``molto importante anche nel videogioco'':
il processo decisionale del giocatore, dallo stimolo all'azione, misurato in
millisecondi: circa \textbf{100 ms per percepire}, \textbf{70 ms per
elaborare}, \textbf{70 ms per il movimento} (premere il pulsante)\footnote{I
valori corrispondono ai tempi tipici del \emph{Model Human Processor} di Card,
Moran e Newell (1983), un modello classico dell'interazione uomo-macchina, più
che a uno studio specifico sui videogiocatori.}.

\subsection{Che cosa fare davanti a una consegna}

Quindi, quando si riceve una consegna, una mail, o il docente dice ``dovete fare
questo per la data X'':
\begin{enumerate}
  \item \textbf{leggere tutto};
  \item \textbf{contare fino a dieci};
  \item \textbf{rileggere con un obiettivo diverso}: fare le quattro colonne
        (richiesta, livelli, perimetro, implicito);
  \item \textbf{elaborare una risposta};
  \item \textbf{rileggere prima di inviare}. ``Rileggete le mail, vi prego.
        \emph{Salve Danilo}, \emph{buongiorno professor Danilo}, \emph{salve
        professor Furetto}\dots: rileggete le mail''.
\end{enumerate}

\subsection{Come si scrive una mail a un docente}

Una mail a un docente è un \textbf{testo operativo}, e il registro adeguato è
\textbf{formale ma non burocratico}. Il relatore spiega come scrivere le mail in
AIV ``perché sono stufo'': fino all'anno prima i docenti correggevano le mail
degli studenti, ma molti non hanno il tempo di farlo (tabella~\ref{pt:tab:mail}).

\begin{table}[htbp]
\centering
\small
\begin{tabular}{@{}lp{9.8cm}@{}}
\toprule
\textbf{Parte} & \textbf{Regola} \\
\midrule
Oggetto & chiaro e specifico; mai vuoto (``le mail con l'oggetto vuoto le
cestino'') \\
Apertura & può essere informale (``Ciao Daniele'': il relatore vuole il tu),
ma rispettosa e diretta \\
Presentazione & se è la prima volta, dire chi si è \\
Corpo & la richiesta, possibilmente una per paragrafo; niente burocratese \\
Chiusura & ``cordiali saluti'' e sempre nome e cognome \\
Firma & impostarla una volta in Gmail: nome e cognome, telefono, mail; per un
docente non proprio, anche corso e anno \\
Tono & informale non vuol dire slang: ``non sono il tuo bro''. Se il docente
risponde in modo ancora più informale si può alleggerire, ma resta come un
capo: ``al vostro capo non scriverete mai \emph{ehi Johnny}'' \\
\bottomrule
\end{tabular}
\caption{Le regole per le mail ai docenti di AIV date dal relatore.}
\label{pt:tab:mail}
\end{table}

Il meme sulla slide è contro il burocratese: nella game industry ``il tempo è
denaro''. Non serve ``Caro professore, le invio la presente per ricordarle il
nostro appuntamento delle ore 14''; serve ``Ciao Daniele, ti ricordi
l'appuntamento delle due e mezza? Ti aspetto nel link''. Fine.

% ---------------------------------------------------------------------
\section{Terzo esercizio: riscrivere una mail}
\label{pt:sec:es3}

L'ultimo esercizio, sempre a coppie, in dieci minuti: riscrivere una mail reale
ricevuta l'anno prima da un team (alla slide 27, il cui link il relatore mette
in chat), applicando quanto visto, ``almeno un paio di cose'': ``ci sono un
fracco di errori nascosti''. La mail era stata inviata anche ai direttori di
corso: ``quindi una mail importante''. Il relatore ha cambiato solo le
indicazioni tecniche, per non far capire chi fosse il team e quale il progetto;
e precisa poi che metà degli errori erano del team reale, metà li ha aggiunti
lui.

\subsection{Gli errori trovati}

\begin{itemize}
  \item \textbf{L'oggetto}: ``Aggiornamento''. Di che cosa?
  \item \textbf{Nessun riferimento} a chi scrive, a quale team, a quale
        progetto. Se è l'ennesima mail non serve un'introduzione enorme, ma
        basta una riga: ``Aggiornamento team XY. Buongiorno a tutti, ecco gli
        aggiornamenti del team''. Non serve nemmeno ``con la presente vi
        elencherò in ordine cronologico'': troppo lungo.
  \item \textbf{La review}: la mail dice ``venerdì 14 novembre''\footnote{Il
        venerdì 14 novembre cade nel 2025, l'anno della mail reale.} alle 14, ma
        non dice \textbf{dove}: ``c'è un posto nell'etere dove fare questa
        riunione?''.
  \item \textbf{Il modo in cui si pongono}, che il relatore considera un problema
        ``enorme'': la mail fissa la review e chiede di avvisare ``entro lunedì''
        se non si può. Ma si sta scrivendo ai direttori di corso: si può
        proporre (``vorremmo proporre una review per venerdì 14 novembre alle
        14: ci sareste, o ci date altri slot?''), non ordinare. ``Chi sei? Non
        sei il mio capo''.
  \item \textbf{L'impaginazione} (lo nota Filippo): la mail voleva comunicare
        due cose, i progressi e i prossimi passi, ma ``visivamente non si
        intuisce'' (lui l'ha capito solo leggendo due volte). Servivano
        macro paragrafi e grassetti: programmazione, game design, animazione,
        narrativa. Giordano, nella sua versione, ha messo un piccolo titolo
        ``Progressi'' e poi ``Milestone'', con i reparti in grassetto.
  \item \textbf{La scansione temporale} (lo nota Giordano): cose fatte, cose in
        corso e cose da fare sono mescolate, e la lettura si incasina.
  \item \textbf{Termini confusi}: ``framework fatto, pretesto narrativo''.
        ``Pretesto narrativo si capisce, ma lo capisco io; e si capisce che
        non è scritto da un narrative, perché quella cosa in narrativa non
        esiste: si può dire \emph{concept}, si può dire \emph{pitch}, ci sono
        termini tecnici; quello non è un termine, è una parola''. E ``da
        riprendere col team'': che cosa significa?
\end{itemize}

% ---------------------------------------------------------------------
\section{Le domande finali}
\label{pt:sec:domande}

Dopo la bibliografia, il relatore resta ``dieci minuti, un quarto d'ora'' per
le domande e per sapere se la lezione è servita: ``molte di queste cose magari
voi le sapete, le date per scontate, ma sappiate che avrete compagni che non
hanno queste conoscenze di base''. È anche l'ultima occasione, per un po', in
cui gli studenti che non fanno narrative design lo vedranno.

\subsection{Faremo comprensione del testo anche in inglese?}

Una studentessa di narrative chiede se si farà comprensione del testo anche in
inglese. Nel corso di narrative non si fa inglese di base, ma c'è un modulo di
\textbf{writing} in cui si faranno cose simili, non dal punto di vista delle
soft skill ma da quello del designer e del giocatore: gli idiomi, come si
scrive una frase in inglese, la struttura sintattica, la punteggiatura. Non è
grammatica inglese dalle basi (``non c'è \emph{the pen is on the table}''):
bisogna già sapere la grammatica di base e scrivere un minimo; ma si danno molti
consigli su come scrivere e farsi capire.

E si lavorerà molto sull'\textbf{intento drammatico} (o intento narrativo):
quando qualcosa ha un obiettivo, il giocatore deve portarsi a casa qualcosa, ma
quell'obiettivo non è in un punto preciso, è nell'insieme. La battuta classica:
il personaggio scende da qualche parte e dice ``oddio, mi devo sbrigare a
trovare questa chiave prima che arrivino i nemici''. L'intento è dire al
giocatore ``attento, stanno arrivando i nemici''; ma c'è di mezzo la narrativa,
quindi non lo si può dire così, lo si deve dire dal punto di vista del
protagonista. ``Nei film esiste, ma nei giochi è ancora peggio''.

\subsection{Puntare al tripla A o a qualcosa di più stabile?}

Filippo fa una domanda ``complessa e arzigogolata'', nata ``questa notte''
(``non voglio risultare né arrogante né criticare''). Gli è sembrato che la
formazione di AIV miri a lavorare soprattutto nelle grandi aziende, nelle
\emph{corporate}. Ma negli ultimi anni c'è una forte crisi del tripla A, anche
di stabilità lavorativa: ha senso, per chi esce e si affaccia per la prima volta
al lavoro, avere quello come obiettivo, o è meglio un ruolo meno remunerativo,
più piccolo, ma più stabile? Se una persona con vent'anni di esperienza viene
licenziata perché chiude il team, forse trova un'altra sistemazione; ma un
giovane senza esperienze, se dopo un anno chiude lo studio?

``Bellissima domanda''. Il relatore risponde ``a step''.

\paragraph{Perché prepararsi per i tripla A.} Anzitutto AIV prepara anche
all'indie: al secondo e al terzo anno, nei progetti, gli studenti saranno in
team di circa quindici persone, ``rimboccarsi le maniche e fare un po' di
tutto'': quella è l'esperienza indie. Ma la domanda è corretta: perché
prepararsi al tripla A? Perché per entrare in Riot servono \textbf{almeno nove
round di colloqui}\footnote{I numeri di round di colloquio citati per Riot,
Larian e CD Projekt sono riportati come li dice il relatore e non sono
verificabili; variano comunque per ruolo e periodo.}: se si è preparati a quelli
e poi si fa un colloquio per un indie, che ne fa due, si è più preparati dello
standard. Riot ne fa nove perché vuole tenersi il lavoratore almeno un anno, ed
essere certa che non abbia problemi due mesi dopo; Larian ne fa quattro, CD
Projekt tre. Di solito c'è un round con le risorse umane, uno tecnico, round
intermedi, e l'ultimo è il \textbf{feeling round}: si lavora con il team, si fa
un meeting con loro per vedere se c'è intesa (``è raro essere scartati, ma
possibile''). Le conoscenze non cambiano: un \emph{core loop} si progetta con la
stessa teoria in una grande o in una piccola azienda. Nella grande bisogna però
stare attenti ad altre dinamiche: c'è una gerarchia più forte, e c'è la
\textbf{HiPPO}, \emph{highest paid person's opinion}: in un brainstorming con il
capo, anche senza volerlo, si dà più ragione a lui perché è il più pagato.

\paragraph{Le persone intraprendenti.} Seconda risposta: prepararsi per una
realtà così complessa crea persone \textbf{intraprendenti}. ``Non c'è lavoro?
Te lo crei. Io me lo sono creato''. A gennaio 2018 il relatore aveva mandato in
un anno, in Italia, 2.222 curriculum\footnote{Così nella trascrizione; il numero
potrebbe essere storpiato.}: un colloquio, zero proposte. Con numeri del genere,
che molti della sua età e anche molti giovani hanno avuto, ``ti rendi conto che
il problema non è il candidato, è il sistema''. E allora ci si fa furbi: il
\textbf{networking} è fondamentale, conoscere le persone, creare relazioni; e
\textbf{fare giochi}, per dimostrare di saperli fare (per questo si fa il
portfolio).

\paragraph{L'ageismo, e i junior.} Sulla questione delle assunzioni c'è in
effetti un problema di \textbf{ageismo} nella game industry: molte persone con
vent'anni di esperienza non vengono riassunte perché considerate ``vecchie'',
non tanto perché si pensi che non sappiano ripartire, ma perché si preferisce
qualcuno di più giovane, che si paga meno. Quindi, paradossalmente, per un
giovane è più semplice fare un anno di esperienza, uscire e ritrovare lavoro:
ha più duttilità, e ``magari non ha famiglia'', non deve ricollocare tutta la
famiglia. I senior sono furbi nei contratti, con molte clausole; e ultimamente
vengono assunti con l'idea che ``tanto ce li ha l'IA che fa il lavoro dei
junior''. A un certo punto si noterà che certi lavori devono farli gli esseri
umani, e si tornerà ad assumere junior: ``il problema delle molte aziende che
adesso non assumono junior sarà questo''. Il relatore è un ``grande
propositore'' dei junior, che hanno il vantaggio di essere
\textbf{incoscienti}: dire ``mi è venuta un'altra idea, perché non possiamo
farlo in un altro modo?''. Una mente che impara bene è più duttile, e vede
soluzioni che un senior, già indirizzato, non vede (``abbiamo sempre fatto così,
continueremo a fare così''). Quindi: ``rompete un po' le uova nel paniere, con
gentilezza''.

In America, in alcuni contesti, l'ambizione è vista positivamente: nelle
\emph{writers' room} ci si aspetta che i junior facciano domande, e se non le
fanno vengono licenziati. Lo stesso nelle aziende high tech: il fratello del
relatore, ingegnere microelettronico in un'azienda internazionale, gli ha
raccontato di un'azienda americana e britannica che ha \textbf{eliminato i
titoli dai contratti} (junior, mid, quadro, senior), perché si era accorta che
chi aveva scritto ``junior'' nei meeting non parlava, mentre l'azienda voleva
che i junior dicessero anche cose sbagliate. ``È un'azienda illuminata'', ma
multimilionaria: ``ci sarà un motivo per cui va avanti, rispetto per esempio
alla Volkswagen, che ha licenziato 100.000 persone ed è stata tolta dalla
quotazione in borsa''\footnote{Inesatto: Volkswagen ha annunciato a fine 2024
un piano di riduzione di circa 35.000 posti di lavoro in Germania entro il
2030, non 100.000 licenziamenti, e non è stata ritirata dalla borsa.}. Il mondo
è pieno di persone convinte che essere brutali, ``qualcuno deve fare dei
tagli'', sia l'unico modo di andare avanti; il modo vero è essere
\textbf{furbi e intraprendenti}, ``non criminali''.

\paragraph{Non svalutatevi.} Due cose ``fuori dai denti''. Primo: \textbf{non
svalutatevi}. ``Non voglio più sentirvi dire: accetto il lavoro perché pagano un
po' meno, però è più tranquillo''. Bisogna essere pagati il giusto; si può dire
``faccio meno ore e quindi ho uno stipendio inferiore, perché voglio fare anche
altro''. Ci sono lavori nella game industry pagati molto bene: quando il
relatore lavorava in un MMO prendeva 3.500 euro al mese, a partita IVA (con una
pressione fiscale bassa, ``vedevo una gran fetta di quei soldi''), per 28 ore a
settimana. In Italia lo stipendio iniziale più basso è intorno ai 1.500-1.800
euro. ``Poi, se siete disperati, ovviamente accettate; ma non buttatevi giù''.

\paragraph{La passività.} Secondo: si entra nell'industria solo
\textbf{lottando}. Il modo peggiore è uscire da AIV, sparire, non tenere i
contatti con i docenti e ripresentarsi due anni dopo: ``non ho ancora trovato
lavoro''. ``Hai fatto quello che ti ho detto? Hai tessuto connessioni con
l'industria? Hai provato un anno all'estero?'' ``No, sono stato in camera,
mandavo ogni tanto qualche candidatura''. ``Questa è passività'', e vale in
tutti i settori. Altro conto è chi, come il relatore nel 2018, le ha provate
tutte senza trovare nulla: a lui è rimasto il \textbf{freelancing}, e si è
scoperto perfetto per farlo, ``forse perché non avevo peli sulla lingua, e in
azienda bisogna essere un po' diplomatici''. Non è colpevolizzare le persone o
dire che ``questa generazione non sa fare niente'': è che c'è davvero chi pensa
che un curriculum al mese basti. Non basta. E i social sono importantissimi, ma
``a volte portano più sponsor che datori di lavoro'': bisogna stare attenti a ciò
che si semina.

\paragraph{I feedback e gli stage.} Filippo precisa che la sua domanda non
criticava la preparazione, ma riguardava gli obiettivi personali. ``Le critiche
le accetto'', risponde il relatore, e aggiunge una cosa importante: a metà e a
fine anno arriveranno dei \textbf{questionari anonimi} per dare feedback ai
docenti. ``Odio il fatto che su dieci studenti me li compilino in cinque, e poi
scopro dal gossip che gli studenti non si trovano bene: mi fa imbestialire''.
Anche ``porca trota, questa cosa non va bene'': si cercherà di sistemarla, nei
limiti tecnici. E ancora, ``non mi nascondo dietro un dito'': AIV non può
assicurare gli \textbf{stage}. Molte accademie in Italia lo promettono, ma gli
stage non esistono più: le scuole li chiedono alle aziende e non ci sono. Li
ottengono solo gli \textbf{ITS}, che sono enti governativi e hanno accordi con
aziende partner. Un'accademia privata non può farlo, e AIV non ha né i
prerequisiti né la voglia di diventare un'università. ``Nessuno può darvi la
garanzia'': spesso chi fa lo stage non viene nemmeno assunto. Conta fare
esperienza, fare giochi, tessere relazioni, anche in fiera: alla Milan Games
Week e Cartoomics a novembre si può parlare con qualcuno e scoprire chi è. Amici
del relatore hanno chiacchierato con ``un tizio giapponese fichissimo'' che poi
ha dato loro il biglietto da visita: era il vicepresidente di Sony. E, al
contrario, attenzione a lamentarsi o a fare gossip in fiera: se si parla male di
Xbox e al tavolo accanto c'è il direttore del marketing di Xbox, ``figura di
merda''.

\subsection{Alcuni corsi trovano lavoro più facilmente di altri?}

Federico chiede se ci siano corsi con una percentuale di occupazione più alta.
``È ciclico''. Ogni anno AIV riceve da aziende richieste di contatti di
studenti, e ogni anno la disciplina cambia: l'anno scorso, a giugno, erano
richiesti tanti programmatori; l'anno prima tanti 3D artist; due anni fa i
game designer\footnote{La trascrizione è contraddittoria: il relatore dice
prima ``l'anno scorso è l'anno dei programmatori'' e poi ``l'anno scorso, a
giugno, ci chiedevano tanti 3D artist''. Si è seguita la prima affermazione.}.
C'è una differenza fra Italia ed estero: per il narrative design, ``aspettatevi
esperienze estere, punto''. Anche in Italia si lavora spesso in inglese, e molti
studi italiani lavorano per Xbox o PlayStation senza che si sappia, con accordi
dietro le quinte (``tu mi fai il development di questa parte di gioco''). Oggi
alcuni tipi di programmatori, non tutti, sono meno richiesti per ``la follia
dell'IA'', ma l'anno prossimo potrebbero tornare a esserlo.

In generale, però, in questi anni c'è un \textbf{calo delle offerte di lavoro su
tutta la linea}: è più difficile entrare, quindi è ancora più importante
prepararsi bene. ``Non vi voglio far iniziare il primo anno con un treno in
faccia, ma meglio saperlo adesso che fra sei mesi''. E vale per tutti i
percorsi: amici delle scuole di fumetto dicono le stesse cose, ``anzi peggio'';
gli architetti; gli informatici, che prima del COVID hanno vissuto una ``moria''
per saturazione (``fai tre anni di informatica e trovi lavoro'', finché non ci
sono stati centomila informatici). Il modo per emergere è fare un buon lavoro,
lasciare una buona impressione, tessere buone conoscenze. Il relatore, da
freelance, non ha un sito curatissimo: i clienti arrivano tutti dal
\textbf{passaparola}, amici di amici; quando gli manca un lavoro extra, ``il
giorno dopo mi chiama uno: ho questo problema, me lo risolvi?''. E crederci:
essere \emph{confident}, alla fine dei tre anni, in ciò che si può dare.

\subsection{Alcune professioni sono pagate meno?}

Un'ultima domanda: è vero che alcune professioni del game dev sono pagate meno
di altre? Sì. Il programmatore è molto ben pagato, soprattutto in America e in
Canada; il game designer lievemente meno; storicamente i meno pagati, ``ed è
brutto'', sono gli artisti. Ma dipende dall'area: in Europa game designer e
programmatori sono più o meno lì; forse i network programmer o chi fa engine
vengono pagati di più, perché spesso la paga dipende dall'effort richiesto. E
tutti i dati online sono \textbf{statistici}: se la media di un programmatore in
America è 100.000 all'anno, c'è chi ne prende 70.000 e chi 160.000 (si vedano
anche gli strumenti del capitolo~\ref{p4:cap:animation}). In Italia ``tutto è
meno riconosciuto'': lo stipendio iniziale, se si è fortunati, è fra i 1.400 e i
1.800 euro, che all'estero a volte è un part time. Nel nord Europa la manodopera
costa tantissimo, e molte aziende spostano la base operativa nell'Est Europa o
nel Sud-est asiatico.

Bisogna considerare anche il costo della vita, osserva Federico. In America la
pressione fiscale è più bassa che in Italia, dove per un dipendente è ``del
55\%, abominevole''\footnote{Il dato è un'approssimazione per eccesso: il cuneo
fiscale italiano su un lavoratore dipendente medio, comprensivo dei contributi a
carico del datore, è stimato dall'OCSE intorno al 45-47\%.}; ma l'assicurazione
medica si paga da sé, e non pagare le tasse è un reato penale. In Svezia,
Norvegia e Danimarca lo stipendio medio è molto più alto, ma anche il costo
della vita (in Finlandia ora c'è una crisi ``spaurosa''). Il \textbf{full remote}
è l'ideale: se offrono una posizione da Helsinki in full remote dall'Italia,
``dite di sì''. In Olanda gli stipendi sono altissimi, ma tutto costa uno
sproposito: un cliente olandese del relatore fa pagare un e-commerce 15.000
euro (il relatore ne chiede 2.500-3.000), ma un caffè gli costa 15 euro.

E poi c'è lo \textbf{stile di vita}: ``l'Italia ancora non ci hanno battuto, il
cibo, lo stile di vita''. Soprattutto in una città come Torino, con molte
aziende videoludiche, dove si gira in bicicletta: ``si fa crunch, ma lo stile di
vita è easy''. Paese che vai, cultura che trovi: in America ``è lavoro, lavoro,
lavoro''; la Danimarca ha una buona cultura del lavoro; in Francia sono molto
ligi, ``alle cinque finisci di lavorare, punto, e se mi chiama il capo non
rispondo'', grazie a un sindacalismo fortissimo: quando si voleva chiudere
Arkane, il ministero avrebbe mandato consulenti del lavoro per tenerla
aperta\footnote{Non verificabile nella forma riportata: nel 2024 Microsoft ha
chiuso Arkane Austin, negli Stati Uniti, mentre la sede francese di Arkane
(Lione) è rimasta aperta.}. In America, se si vuole chiudere si chiude; in
Italia, come in Francia, si cerca sempre di salvare: a volte è eccessivo, a
volte è importante.

\subsection{I saluti}

``Grazie mille, è stato un piacere''. Ai narrative designer: ``ci vediamo domani
mattina''; agli altri: ``mi vedrete soprattutto al secondo anno, in control
room, a fare il mastino da guardia''. E un ultimo appello: ``vi prego, fate dei
bei giochi, perché abbiamo bisogno di dimostrare come italiani che siamo
creativi e sappiamo fare un fracco di bellissimi giochi. Ci serve un
\emph{Expedition 33} nei prossimi anni. Non vi voglio dare un peso sulle spalle,
ma ve lo sto dando: accollatevelo''.
